% 10-appendix.tex
\section*{附录：符号与核心清单}

\paragraph{符号。}
\texttt{path:line} 表示相对于存储库根目录的文件；是该文件的第 49 行，并且 \texttt{path:line--line} 是
包含范围。代码清单是逐字摘录，最多包含空格折叠；标题给出了上下文。表格列举了规范值（关键数字、工件名称、修剪轴），而不是
说明性样本。

\paragraph{中央实体。}
对于想要通过核心的最短路径的读者来说，
中心定义是：调度表（\texttt{array<atomic<KernelFunction>,13>}），
键枚举、张量对象模型 (\texttt{Tensor}/\texttt{TensorImpl}/\texttt{Storage})、
自动梯度图 (\texttt{Node}/\texttt{Edge}/\texttt{GraphTask}/\texttt{InputBuffer}/\texttt{Engine}
与 \texttt{SavedVariable} 防护和 \texttt{VariableVersion} 计数器），
异常模式、梯度/推理 TLS 模式、分析器的
\texttt{OpRecord}/\texttt{GpuTimerPair} 堆栈，偏移量 9 处的 vmap 键，以及
代码生成背后的两个 YAML 注册表。每一章都以一个图开始
将其实体放置在四柱地图上。